\chapter{Transformations}\label{geometry:ch:transformations}

\begin{center}\itshape
Given a manifoldness and a group of transformations of the same; to investigate the configurations belonging to the
manifoldness with regard to such properties as are not altered by the transformations of the group.
\upshape(Felix Klein, the Erlangen programme, 1872, in the translation of 1893)
\end{center}

A reflection changes a picture but keeps every distance. A stretch can change angles while keeping parallel lines parallel. Which conclusions survive the transformations we allow? Klein's 1872 Erlangen programme proposed this as a way to organize geometry. A larger transformation group preserves fewer properties, so it asks different questions about the same figures.

The idea reaches familiar patterns. Fedorov classified the seventeen wallpaper groups in 1891; P\'olya drew them again in 1924. Escher, who studied the Alhambra's tilings during visits in 1922 and 1936, used P\'olya's work in his repeating designs. We will prove the smaller seven-type classification for friezes, where repetition occurs along a strip. The wallpaper and space-group counts will remain quoted results.

The tools are groups and group actions from Volume~3, Chapters~1--4, together with Euclidean vectors and matrices. Chapters~7 and~9 have already used transformation groups for matrix manifolds and hyperbolic geometry. Here the plane makes the same method visible through mirrors, rotations and lattice vectors.

\begin{think}
Imagine a repeating row of footprints. Walking alternates left and right feet; hopping repeats one foot; jumping can leave both together. Which patterns have a mirror symmetry? Which agree with themselves only after a reflection followed by a shift?
\end{think}

\section{Isometries}

An \emph{isometry} of $\R^n$ is a map preserving all distances: $|f(x)-f(y)|=|x-y|$. Surjectivity is not included in this definition; in finite-dimensional Euclidean space it follows.

\begin{theorem}[isometries are affine]\label{geometry:thm:isometryaffine}
Every isometry of $\R^n$ is $f(x)=Ax+b$ for an orthogonal matrix $A$ and vector $b$. Hence it is bijective, and all isometries form a group under composition.
\end{theorem}
\begin{proof}
Subtract $b=f(0)$ and call the resulting map $g$. It fixes zero, so $|g(x)|=|x|$. Expanding squared distances gives $\langle g(x),g(y)\rangle=\langle x,y\rangle$. In positive dimension the vectors $g(e_i)$ form an orthonormal basis. The coordinate of $g(x)$ along $g(e_i)$ is $\langle g(x),g(e_i)\rangle=x_i$, so
\[
g(x)=\sum_i x_i g(e_i).
\]
Thus $g$ is linear with orthogonal matrix $A$. Its inverse is $A^\top$, and $f^{-1}(x)=A^\top(x-b)$ preserves distances. Compositions preserve distances as well. In dimension zero there is only the identity map on a point, and the same formula uses the empty matrix and vector.
\end{proof}

Identify the plane with $\C$. The first column of an orthogonal matrix is $(\cos\theta,\sin\theta)$, and its second is either $(-\sin\theta,\cos\theta)$ or its negative. These give a rotation or a reflection through the origin. Every plane isometry consequently has one of the forms
\begin{equation}\label{geometry:eq:isomC}
\begin{gathered}
z\longmapsto e^{i\theta}z+b\quad\text{(direct)},\\
z\longmapsto e^{i\theta}\bar z+b\quad\text{(indirect)}.
\end{gathered}
\end{equation}
The words mean orientation-preserving and orientation-reversing, respectively.

\begin{theorem}[three reflections suffice]\label{geometry:thm:threereflections}
The images of three noncollinear points determine a plane isometry. Every plane isometry is a product of at most three reflections in lines.
\end{theorem}
\begin{proof}
If $P,Q$ have the same distances to $A,B,C$, subtracting their squared-distance equations shows that $P-Q$ is perpendicular to both $B-A$ and $C-A$. Noncollinearity makes these a basis, so $P=Q$. Thus agreement at $A,B,C$ forces agreement everywhere.

To realize a given $f$, first reflect $A$ to $f(A)$ across their perpendicular bisector, unless they already agree, in which case do nothing. Call this map $s_1$. Next reflect $s_1(B)$ to $f(B)$, again omitting the reflection if unnecessary. Their perpendicular bisector contains $f(A)$, since the two points have the same distance from it. Thus this second map $s_2$ keeps the first image fixed. Finally the two candidate images $s_2s_1(C)$ and $f(C)$ have equal distances from $f(A),f(B)$. If they differ, their perpendicular bisector contains both of those points, and reflection in that line matches the last image while fixing the first two. The resulting product agrees with $f$ at all three points, hence is $f$. An omitted step contributes no reflection, including all three steps for the identity.
\end{proof}

\begin{figure}[ht]
\centering
\fig{geometry-ch02-tworeflections}
\caption{A reflection in $m_1$ followed by a reflection in $m_2$ is the rotation about their intersection by twice the angle
$\varphi$ from $m_1$ to $m_2$ (Exercise~\ref{geometry:ex:chasles}).}
\label{geometry:fig:tworeflections}
\end{figure}

Two intersecting mirrors give twice their directed angle as a rotation; two parallel mirrors give twice their perpendicular separation as a translation. Exercise~\ref{geometry:ex:chasles} checks the formulas, including coincident mirrors. The following classification can be read directly from the complex expression, without assuming those formulas first.

\begin{theorem}[classification of plane isometries]\label{geometry:thm:chasles}
Every plane isometry belongs to exactly one of five classes: the identity; a nonzero translation; a rotation about a point with nonzero angle modulo $2\pi$; a reflection in a line; or a \emph{glide reflection}, a reflection followed by a nonzero translation along its mirror line.
\end{theorem}
\begin{proof}
For a direct map $f(z)=e^{i\theta}z+b$, the case $e^{i\theta}=1$ is a translation, the identity exactly when $b=0$. Otherwise $c=b/(1-e^{i\theta})$ is its unique fixed point and $f(z)-c=e^{i\theta}(z-c)$, a rotation about $c$.

For an indirect map, use $r(z)=e^{i\theta/2}z$ to obtain
\[
r^{-1}fr(z)=\bar z+b',\qquad b'=e^{-i\theta/2}b=t+iu.
\]
In real coordinates this is $(x,y)\mapsto(x+t,u-y)$: reflect in $y=u/2$, then shift along that line by $t$. It is a reflection for $t=0$ and a genuine glide for $t\ne0$. Rotating coordinates back gives the same description for $f$. The classes are distinct by orientation and fixed sets: the direct fixed sets are the whole plane, empty, or one point; the indirect ones are a line or empty.
\end{proof}

A \emph{similarity} multiplies all distances by one constant $k>0$. Dividing its target coordinates by $k$ gives an isometry, so plane similarities are $az+b$ and $a\bar z+b$ with $a\ne0$. For a direct similarity with $a\ne1$, the fixed point is $b/(1-a)$, and about that point the map rotates and scales. This is a \emph{spiral similarity}. The case $a=1$ is a translation, possibly the identity.

\section{Klein's Erlangen programme}

The following inclusions organize several familiar properties. For the last row, extend affine maps to the projective plane; Appendix~C will define that plane and its cross-ratio. This row is a preview, not a tool needed in the proofs here.

\medskip
\begin{center}
\begin{tabular}{@{}p{0.44\linewidth}p{0.5\linewidth}@{}}
\hline
group & preserved properties\\
\hline
isometries $x\mapsto Ax+b$, $A$ orthogonal & lengths, angles, areas\\
similarities $x\mapsto kAx+b$, $k>0$ & angles, length ratios, circles\\
affine maps $x\mapsto Mx+b$, $M$ invertible & lines, parallelism, ratios along a line, area ratios\\
projective maps (Appendix~C) & lines, cross-ratio\\
\hline
\end{tabular}
\end{center}
\medskip

An affine map sends $p+tu$ to $Mp+b+tMu$, preserving lines, parallelism and the line parameter $t$. Its area multiplier is $|\det M|$, so ratios of areas survive too. Thus a statement about medians or side-division ratios can often be checked on one convenient triangle.

\begin{proposition}\label{geometry:prop:affinetriangles}
For any two nondegenerate triangles $ABC$ and $A'B'C'$, exactly one invertible affine map sends their respective vertices to each other.
\end{proposition}
\begin{proof}
The map must be $x\mapsto A'+M(x-A)$, with $M(B-A)=B'-A'$ and $M(C-A)=C'-A'$. The two source vectors form a basis, so these equations determine one linear map. The target vectors also form a basis, making it invertible.
\end{proof}

An affine triangle problem can therefore be moved to $(0,0),(1,0),(0,1)$, as in the Routh exercise below. An ellipse, viewed as an invertible affine image of a circle, supplies another example. Parallel chords and their midpoints move to parallel chords and their midpoints, and the circle's centre moves to the ellipse's centre. Their midpoint-line property is consequently affine even though perpendicularity is not.

\section{Finite symmetry groups}

The \emph{symmetry group} of a figure $X$ consists of all isometries carrying $X$ onto itself. A regular $n$-gon, $n\ge3$, has the dihedral group $D_n$ of order $2n$: $n$ rotations and $n$ reflections. A pinwheel may keep the rotations but lose the mirrors, leaving a cyclic group $C_n$. Leonardo studied such rotationally symmetric designs; Weyl associated his name with the classification below.

\begin{theorem}[Leonardo's theorem]\label{geometry:thm:leonardo}
A finite group of plane isometries has a common fixed point. About that point it is either a cyclic group $C_n$ of rotations or a dihedral group $D_n$ with $n$ rotations and $n$ reflections, for some $n\ge1$.
\end{theorem}
\begin{proof}
For any $p$, average its orbit with multiplicities:
\[
c=\frac1{|G|}\sum_{g\in G}g(p).
\]
An isometry is affine and preserves such averages. For $h\in G$, multiplication permutes the summands, so $h(c)=c$. Move $c$ to zero. Every group element is now orthogonal, hence a rotation or reflection through zero.

Let $G^+$ be the rotation subgroup, the kernel of orientation. If it is trivial, put $n=1$. Otherwise choose its least positive rotation angle $\theta\in(0,2\pi)$. Dividing any other angle by $\theta$ and subtracting a multiple leaves a rotation with angle in $[0,\theta)$, so its remainder must be zero. Applying the same argument to $2\pi$ gives $\theta=2\pi/n$. Thus $G^+=C_n$. If there are no reflections, $G=G^+$. If there is a reflection $s$, the other orientation coset is $sG^+$, giving $n$ reflections and the relation $srs=r^{-1}$ for a generating rotation $r$. This is $D_n$. For $n=1$ the two possibilities are the trivial group and a group with a single reflection; neither requires interpreting a one-vertex figure as a regular polygon.
\end{proof}

\section{Friezes}

\begin{definition}\label{geometry:def:frieze}
A \emph{frieze group} is a group of plane isometries preserving the real axis whose translation subgroup consists exactly of integer multiples of one nonzero translation along it. Scale the unit so this subgroup is $T=\{z\mapsto z+k:k\in\Z\}$.
\end{definition}

An isometry preserving the axis has the real-coordinate form
\[
(x,y)\longmapsto(\varepsilon x+t,\eta y),\qquad \varepsilon,\eta\in\{1,-1\}.
\]
The first sign tells whether it reverses direction along the axis; the second tells whether it exchanges the two sides. Their pair defines a homomorphism $\sigma:G\to V=\{1,-1\}^2$ with kernel $T$. The four kinds are translations $z+t$, half-turns $-z+t$, vertical mirrors $-\bar z+t$, and horizontal mirrors or glides $\bar z+t$, respectively.

Two groups have the same \emph{type} if a similarity preserving the real axis conjugates one to the other. Changing the origin along the axis can move half-turn centres and vertical mirrors. It cannot turn a genuine horizontal glide into a horizontal mirror.

\begin{theorem}[the seven friezes]\label{geometry:thm:friezes}
There are exactly seven frieze types, generated by $T$ and the following extra maps:
\begin{center}
\begin{tabular}{@{}lll@{}}
\hline
$\sigma(G)$ & extra generators & name\\
\hline
$\{1\}$ & none & hop\\
$\langle(+,-)\rangle$ & $\bar z+\frac12$ & step\\
$\langle(-,+)\rangle$ & $-\bar z$ & sidle\\
$\langle(-,-)\rangle$ & $-z$ & spinning hop\\
$V$ & $-\bar z$, $\bar z+\frac12$ & spinning sidle\\
$\langle(+,-)\rangle$ & $\bar z$ & jump\\
$V$ & $-\bar z$, $\bar z$ & spinning jump\\
\hline
\end{tabular}
\end{center}
These names are Conway's. The spinning sidle has vertical mirrors, glides and half-turns; the spinning jump additionally has the horizontal mirror. The other rows have the symmetries described by their generators and translations.
\end{theorem}
\begin{proof}
The sign image is one of the five subgroups of $V$: trivial, one of its three order-two subgroups, or all of $V$. Because the kernel is $T$, each sign has exactly one coset of translations. For $r_s(z)=z+s$, the conjugation $r_s^{-1}fr_s$ changes the parameter $t$ to $t-2s$ in $-z+t$ and $-\bar z+t$, but leaves it unchanged in $\bar z+t$.

The trivial image gives $G=T$, the hop. An image generated by $(-,-)$ has a half-turn $-z+t$; choose $s=t/2$ to centre it at zero. This gives the spinning hop. The same normalization of a vertical mirror gives the sidle when the image is generated by $(-,+)$.

For the image generated by $(+,-)$, choose $h(z)=\bar z+t$. Its square is translation by $2t$, so $2t\in\Z$. Multiplying by an integer translation reduces $t$ to zero or $1/2$. The first case contains the horizontal mirror and is the jump; the second has only horizontal glides and is the step.

For full image $V$, first centre a vertical mirror $s(z)=-\bar z$, then take $h(z)=\bar z+t$ with $t=0$ or $1/2$ as above. These give the spinning jump and spinning sidle. In the latter, $sh(z)=-z-1/2$ is a half-turn centred at $-1/4$. Its translates have centres halfway between the vertical mirrors at the half-integers.

These possibilities really form groups with exactly the required translations. A half-turn or vertical mirror squares to the identity and conjugates integer translation to its inverse. The horizontal generator commutes with translations and has square translation by $2t\in\Z$. In the full-image cases $shs=h^{-1}$, so every word reduces to an integer translation times one of $1,s,h,sh$. Their different signs show that only the first coset contributes translations. The analogous two-coset description proves the same fact for the order-two images. Thus no extra, smaller translation has entered the construction. Finally the sign subgroup and the presence of the horizontal mirror are unchanged by allowed conjugations. They distinguish all seven rows.
\end{proof}

\begin{figure}[ht]
\centering
\fig{geometry-ch02-friezes}
\caption{The seven frieze patterns, each obtained by applying a frieze group to one L-shaped tile. The dotted line is the axis. The
step is the trail of a walker, the jump that of someone jumping with feet together.}
\label{geometry:fig:friezes}
\end{figure}

A hop repeats one footprint; a step combines alternation with a glide; a jump adds a mirror through the middle of the strip. Conway's names make the sign calculation easier to remember. In a physical border or fabric, use the exact geometric symmetry of the design rather than the approximate appearance of a hand-drawn tile. Letter-pattern examples below make the same idealization explicit.

\section{Lattices and the crystallographic restriction}

A \emph{lattice} in $\R^n$ is $\Lambda=\Z v_1\oplus\cdots\oplus\Z v_n$ for a real basis $v_1,\ldots,v_n$. Inverting that basis matrix shows that bounded regions contain only finitely many lattice points: bounded integer coordinates have finitely many choices. In particular nonzero lattice vectors have a positive shortest length. We use full-rank discrete translation groups for periodic plane patterns. Mere repetition is insufficient: a constant pattern, for instance, admits every translation and has no discrete translation lattice. Appendix~H will study lattices further.

\begin{think}
Square tiles accommodate quarter-turns and honeycombs accommodate sixth-turns. Could a full-rank lattice also allow a fifth-turn? An eighth-turn? The trace of its rotation matrix gives a quick test.
\end{think}

\begin{theorem}[the crystallographic restriction]\label{geometry:thm:crystallographic}
If a finite-order rotation $\rho$ preserves a lattice in $\R^2$ or $\R^3$, its order belongs to $\{1,2,3,4,6\}$.
\end{theorem}
\begin{proof}
Write the rotation as $\rho(x)=Ax+b$, allowing its centre to be away from zero. Since $b=\rho(0)\in\Lambda$, differences $\rho(v)-\rho(0)$ show $A\Lambda\subset\Lambda$, and the inverse rotation gives equality. In a lattice basis $A$ therefore has integer entries and integer trace.

In orthonormal coordinates a plane rotation through $\theta$ has trace $2\cos\theta$; a spatial rotation about its axis has trace $1+2\cos\theta$. The latter description was established in Chapter~7. Thus $2\cos\theta$ is an integer in either case, belonging to $\{-2,-1,0,1,2\}$. If the order is $n$, its angle is $2\pi k/n$ with $(k,n)=1$, or zero for the identity. The listed cosine values give exactly orders $1,2,3,4,6$.
\end{proof}

\begin{remark}
Shechtman's 1982 observation, published with Blech, Gratias and Cahn in 1984, revealed quasicrystalline order without translational periodicity. Such diffraction patterns can show rotational symmetries excluded by this theorem. The missing hypothesis is the full-rank translation lattice, so there is no contradiction. Penrose tilings offer a related nonperiodic planar example.
\end{remark}

The \emph{point group} of a plane lattice consists of orthogonal maps preserving it. It is finite: choose a bounded finite set of lattice vectors spanning the plane; the point group permutes the finite set of all lattice vectors up to that length, and its action there determines the map. It contains the half-turn $-I$. Leonardo and the restriction theorem give a cyclic or dihedral group with rotation order $2,4$ or $6$.

\begin{proposition}[the five lattice types]\label{geometry:prop:fivelattices}
A plane lattice's point group is $C_2,D_2,D_4$ or $D_6$. If it contains a reflection $s$, it has a basis either $a,b$, with $a$ on the mirror and $b$ perpendicular to it, or $a,(a+b)/2$ for such perpendicular $a,b$. These give rectangular and centred rectangular descriptions; together with the oblique, square and hexagonal cases they yield the five lattice types in Figure~\ref{geometry:fig:lattices}.
\end{proposition}
\begin{proof}
First exclude cyclic-only point groups of order four or six. Let $v$ be a shortest nonzero lattice vector and let $w$ be its rotation by $90^\circ$ or $60^\circ$, respectively (choose the inverse generating rotation if necessary). They have equal length $\delta$. Any plane vector can be reduced modulo $\Z v+\Z w$ to $av+bw$ with $|a|,|b|\le1/2$. Its squared length is at most
\[
\delta^2\bigl(\tfrac12+\tfrac12|\cos\theta|\bigr)<\delta^2
\quad(\theta=90^\circ\text{ or }60^\circ).
\]
For a lattice vector the remainder must therefore be zero, by shortest length. Thus $v,w$ generate the lattice. Reflection in the line of $v$ fixes $v$ and sends $w$ to $2\cos\theta\,v-w$, an integer combination in either case. It preserves the lattice and supplies the missing reflection, giving $D_4$ or $D_6$.

Now suppose $s$ is a lattice reflection with mirror $\ell$. The vectors $v+sv$ lie on $\ell$ and $v-sv$ on $\ell^\perp$. Each collection has a nonzero member, since otherwise the full lattice would lie on one line. The discrete line subgroups $\Lambda\cap\ell$ and $\Lambda\cap\ell^\perp$ have primitive generators $a,b$, chosen as their shortest vectors in a fixed positive direction. Every $2v$ belongs to $\Lambda_0=\Z a+\Z b$, so $\Lambda_0\subset\Lambda\subset\tfrac12\Lambda_0$.

Modulo $\Lambda_0$, the only possible nonzero classes are $a/2,b/2,(a+b)/2$. The first two cannot occur: subtracting an element of $\Lambda_0$ would give a lattice vector on the corresponding axis shorter than its primitive generator. Therefore the quotient is trivial or generated by $(a+b)/2$. The respective bases are $a,b$ or $a,(a+b)/2$.

With no mirrors the point group is $C_2$, the oblique type. Generic rectangular and centred rectangular lattices have $D_2$; special shapes acquire the extra rotations giving square $D_4$ or hexagonal $D_6$. These exceptional shapes are counted separately rather than counted twice among their rectangular descriptions.
\end{proof}

\begin{figure}[ht]
\centering
\fig{geometry-ch02-lattices}
\caption{The five types of plane lattices, with point groups $C_2$, $D_2$, $D_2$, $D_4$, $D_6$. In the centred rectangular lattice the
rectangle spanned by $a$ and $b$ has a lattice point at its centre.}
\label{geometry:fig:lattices}
\end{figure}

\begin{remark}
A \emph{wallpaper group} is a plane isometry group whose translations form a full-rank lattice. The seventeen-group classification of Fedorov (1891) is quoted here, not proved. Its proof combines lattice point groups with possible mirrors and glides; Armstrong gives the cases, and Conway, Burgiel and Goodman-Strauss explain their orbifold notation. The corresponding count of three-dimensional space groups is $230$, independently obtained by Fedorov and Sch\"onflies in 1891; this too is a quoted classification.
\end{remark}

\section{A gem: the shortest closed billiard path}

For an acute triangle, Fagnano asked in 1775 for the shortest perimeter of a triangle with one vertex on each side. Reflections turn every candidate perimeter into a lower bound by a straight segment. We then need to check that the proposed minimizer actually attains it.

\begin{theorem}[Fagnano]\label{geometry:thm:fagnano}
In an acute triangle $ABC$, the inscribed triangle of minimum perimeter is the orthic triangle, formed by the three altitude feet.
\end{theorem}
\begin{proof}
Take $D$ on $BC$, $F$ on $AB$ and $E$ on $AC$. Reflect $D$ across the two side lines to $D',D''$. Since $DF=D'F$ and $DE=D''E$,
\[
DF+FE+ED=D'F+FE+ED''\ge D'D''=2AD\sin A\ge2h_a\sin A.
\]
Here the angle between the two reflected rays from $A$ is $2A$, and $h_a$ is the altitude from $A$ to $BC$. This lower bound holds for every choice of points, whether or not the segment $D'D''$ meets the two sides in the desired order.

To verify attainment, let $D,F,E$ be the altitude feet from $A,C,B$, all interior to their sides because the triangle is acute. Put $a=BC$, $b=CA$, $c=AB$, and let $R$ be the circumradius. The right triangles give $AF=b\cos A$ and $AE=c\cos A$. By the cosine rule,
\[
EF^2=(b^2+c^2-2bc\cos A)\cos^2 A=a^2\cos^2 A.
\]
Thus $EF=a\cos A$, with the cyclic formulas for the other two orthic sides. Their perimeter is
\[
a\cos A+b\cos B+c\cos C
=R(\sin2A+\sin2B+\sin2C)
=4R\sin A\sin B\sin C=2h_a\sin A.
\]
The trigonometric identity uses $A+B+C=\pi$; for example the sum of the first two double-angle sines is $2\sin C\cos(A-B)$, which combines with $\sin2C$ to give the displayed product. The orthic triangle reaches the lower bound.

Any minimizer must have $AD=h_a$, forcing $D$ to be the unique altitude foot. Equality in the broken-line inequality then forces $D',F,E,D''$ to occur in order on one segment. The orthic construction has just been shown to attain equality, so that segment's unique intersections with $AB,AC$ are its orthic $F,E$. This proves uniqueness as well.
\end{proof}

At each orthic foot the path obeys equal incidence and reflection angles: the unfolded path is straight across the mirror. This is the usual geometric model for billiards and reflected light. A trajectory hitting a triangle vertex has no single side normal there and needs separate rules; the periodic square-billiard exercise excludes such corner collisions.

\begin{beauty}
The same reflection can answer two different questions. In an isometry group it is a generator whose products encode rotations and glides. In a shortest-path problem it replaces a bent path by a straight one, making the triangle inequality sharp. Fagnano's problem needs both halves of that reasoning: a bound for every candidate and an equality check for the actual altitude feet.
\end{beauty}

Ceva's concurrence criterion uses lines and signed ratios, so it survives affine changes of coordinates. The Euler line also involves the circumcircle, tying it to similarities instead. The exercises prove both and show precisely which information each transformation preserves.

\begin{figure}[ht]
\centering
\fig{geometry-ch01-euler}
\caption{The Euler line through the circumcentre, the centroid and the orthocentre, with the centroid twice as far from the orthocentre as from the circumcentre, and the nine-point circle of half the circumradius.}
\label{geometry:fig:euler}
\end{figure}

\begin{figure}[ht]
\centering
\fig{v9-appa-orthic-unfolding}
\caption{At the orthic triangle in this acute equilateral example, reflecting the altitude foot $D$ across $AB,AC$ gives $D',D''$. The straight route $D'\!F\!E\!D''$ has the same length as the closed route $D\!F\!E\!D$. This verifies equality for the minimizing configuration; no such side-intersection claim is needed for arbitrary $D$.}
\label{fig:v9-orthic-unfolding}
\end{figure}

\begin{figure}[ht]
\centering
\fig{v9-appa-fivefold-obstruction}
\caption{If a plane lattice had a $72^\circ$ rotation and a shortest nonzero vector $v$ of length $\delta$, then $v+\rho^2v$ would be a nonzero lattice vector of length $2\delta\cos72^\circ<\delta$. It lies inside the dashed shortest-length circle, a contradiction. The other marked orbit points are hypothetical lattice points under that assumption.}
\label{fig:v9-fivefold-obstruction}
\end{figure}

Translations and mirrors classify a strip through two signs; a lattice adds the arithmetic condition of an integer trace. Affine maps let us move whole families of triangles and ellipses to convenient models, while remembering that lengths and right angles need not survive. Appendix~B changes the transformation again: inversion exchanges lines and circles and keeps angles, bringing configurations that look unrelated into the same picture.

\section{Exercises}

\begin{exercise}[$\star$]\label{geometry:ex:chasles}
Compute a reflection in $m_1$ followed by one in $m_2$. For intersecting lines, prove it is rotation through twice the directed angle from the first mirror to the second. For parallel lines, prove it is translation by twice their perpendicular displacement, including the identity for coincident mirrors.
\end{exercise}
\begin{exercise}[$\star$]\label{geometry:ex:tworotations}
Compose rotation by $\alpha$ about $p$ with rotation by $\beta$ about $q$, in that order. Find the centre when $\alpha+\beta\not\equiv0\pmod{2\pi}$, and the translation vector otherwise. Include coincident centres and the possible identity.
\end{exercise}
\begin{exercise}[$\star$]\label{geometry:ex:dihedral}
For a regular $n$-gon, $n\ge3$, prove that its symmetry group has $2n$ elements, that its rotation subgroup is normal, and that $srs=r^{-1}$ for any reflection $s$ and rotation $r$ in the group.
\end{exercise}
\begin{exercise}[$\star\star$]\label{geometry:ex:conjugacy}
In the full plane isometry group, classify conjugacy within each kind: rotations by angles modulo $2\pi$, translations by their vectors, reflections, and genuine glides. Prove that the respective invariants are angle up to sign, vector length, no further invariant for reflections, and glide-displacement length. Prove sufficiency as well as necessity.
\end{exercise}
\begin{exercise}[$\star$]\label{geometry:ex:affineellipse}
Prove that the midpoints of all chords of an ellipse in one fixed parallel direction lie on a line through its centre. Use an affine circle model and identify the correct diameter direction in that circle before transforming it.
\end{exercise}
\begin{exercise}[$\star$]\label{geometry:ex:friezeid}
Identify the frieze types of
$\dots\mathsf{LLLL}\dots$, $\dots\mathsf{NNNN}\dots$, $\dots\mathsf{VVVV}\dots$, $\dots\mathsf{DDDD}\dots$, $\dots\mathsf{HHHH}\dots$,
$\dots\mathsf{pbpb}\dots$, and $\dots\textsf{VΛVΛ}\dots$.
Use separated, equally spaced ideal geometric glyphs in equal-height cells, not small irregularities of a printed font: $L$ has no mirror or half-turn, $N$ only a half-turn, $V$ only a vertical mirror, $D$ only a horizontal mirror, and $H$ both mirrors. The asymmetric glyphs $p,b$ are horizontal reflections of each other; $\Lambda$ is the horizontal reflection of $V$. For the alternating rows, find the smallest translation and any half-period glide.
\end{exercise}
\begin{exercise}[$\star$]\label{geometry:ex:fivefold}
Without using trace, rule out a fifth-order rotation preserving a plane lattice. Reduce to a $72^\circ$ linear rotation, choose a shortest nonzero lattice vector, and produce a shorter one.
\end{exercise}
\begin{exercise}[$\star\star$]\label{geometry:ex:hexagonal}
Let a plane lattice be preserved by a $60^\circ$ rotation $\rho$ about zero, and choose a shortest nonzero vector $v$. Prove $\Lambda=\Z v+\Z\rho v$ and deduce similarity with the Eisenstein lattice $\Z[\omega]$, where $\omega=e^{2\pi i/3}$.
\end{exercise}
\begin{exercise}[$\star\star$]\label{geometry:ex:cube}
Count the cube's $24$ rotations in two ways: use the image of one face with a marked edge, then list rotations by their axes. Explain why the count includes every rotation and has no repetitions.
\end{exercise}
\begin{exercise}[$\star\star$]\label{geometry:ex:billiards}
A square-billiard trajectory obeys equal reflection angles and never hits a corner. With initial direction proportional to $(q,p)$, where $q>0$, $p\in\Z$ and $(|p|,q)=1$, prove it is periodic in both position and direction and find its shortest positive length. Include the horizontal case, and state the analogous vertical result. Use unfolded reflected squares.
\end{exercise}
\begin{exercise}[$\star$]\label{geometry:ex:glidesquare}
Prove that a genuine glide's square is a nonzero translation. For a prescribed such translation, describe all glides with that square and explain how choosing the axis determines the glide. What happens for the identity translation?
\end{exercise}
\begin{exercise}[$\star\star\star$]\label{geometry:ex:finiteSO3}
Let $G$ be a finite rotation group of $\R^3$ fixing zero, with order $N>1$. Each nonidentity rotation has two poles on $S^2$. Count pairs (nonidentity element, fixed pole) to obtain
\[
2\left(1-\frac1N\right)=\sum_i\left(1-\frac1{n_i}\right),
\]
where $i$ runs over pole orbits and $n_i$ is the stabilizer order. Prove that the orbit orders are one of $(n,n)$, $(2,2,n)$, $(2,3,3)$, $(2,3,4)$ or $(2,3,5)$. Determine the group order from each list.
\end{exercise}
\begin{exercise}[$\star\star$]\label{geometry:ex:ceva}
For a nondegenerate triangle $ABC$, take finite points $D,E,F$ on the lines $BC,CA,AB$, excluding the vertices. Prove that $AD,BE,CF$ are concurrent or parallel exactly when
\[
\frac{BD}{DC}\frac{CE}{EA}\frac{AF}{FB}=1.
\]
Ratios are signed: $D-B=(BD/DC)(C-D)$, with cyclic conventions for the other sides. Include both alternatives in both directions of the proof.
\end{exercise}
\begin{exercise}[$\star$]\label{geometry:ex:cevaparallel}
Give a direct proof of the parallel alternative in Ceva's criterion by choosing coordinates in which the three cevians are vertical. Check why all denominators in the signed-ratio calculation are nonzero.
\end{exercise}
\begin{exercise}[$\star\star$]\label{geometry:ex:eulerline}
Put a nondegenerate triangle's circumcircle on $|z|=1$, with vertices $a,b,c$. Prove that its orthocentre is $h=a+b+c$ and centroid $g=(a+b+c)/3$. Deduce collinearity with the circumcentre $0$ and $|h-g|=2|g|$; when distinct, $g$ lies between the other centres. Prove that all three coincide exactly in the equilateral case, as in Figure~\ref{geometry:fig:euler}.
\end{exercise}
\begin{exercise}[$\star\star$]\label{geometry:ex:routh}
In a triangle of area one, divide the sides in ratios $1:2$ while going consistently around the boundary. Join each vertex to its opposite division point. Prove that the three cevians enclose area $1/7$, keeping track of both the intersection coordinates and the area scale.
\end{exercise}

\begin{exercise}[$\star\star\star$]\label{ex:v9-finite-affine-group}
\emph{A finite affine group has an invariant ruler.} Let $G$ be a finite group of invertible affine maps of $\R^2$. (a) Prove it has a common fixed point $c$. In coordinates centred at $c$, write its matrices as $A_g$ and prove that $Q=\sum_gA_g^\top A_g$ is positive definite with $A_h^\top QA_h=Q$. (b) Construct an invertible real $S$ with $S^\top S=Q$ by a local $2\times2$ calculation, and show that $x\mapsto S(x-c)$ conjugates $G$ to an orthogonal cyclic or dihedral group. (c) Give a finite affine group that is not a Euclidean isometry group in its original coordinates.
\end{exercise}
\begin{exercise}[$\star\star$]\label{ex:v9-glide-wallpaper}
\emph{A periodic pattern with glides but no mirrors.} Let $g(x,y)=(x+1/2,-y)$ and $b(x,y)=(x,y+1)$, and put $a=g^2$. (a) Give unique normal forms for $G=\langle g,b\rangle$ and find its translation lattice and quotient by translations. (b) Find every glide axis and show that no nonidentity element fixes a point. (c) Show that $[0,1/2)\times[0,1)$ contains exactly one representative of each orbit. Explain why this group is not conjugate to the group generated by unit horizontal and vertical translations and a horizontal reflection.
\end{exercise}
\begin{exercise}[$\star\star$]\label{ex:v9-ellipse-parallelogram}
\emph{The largest parallelograms in an ellipse.} Let $E=c+T(S^1)$ for an invertible real $2\times2$ matrix $T$. (a) Find the greatest area of a parallelogram with its four vertices on $E$ and characterize all equality cases. (b) For $T=\operatorname{diag}(2,1)$, compare the diamond with vertices $(\pm2,0),(0,\pm1)$ and the rectangle with vertices $(\pm\sqrt2,\pm1/\sqrt2)$. Explain why the maximizers need not all have right angles in the ellipse's Euclidean metric.
\end{exercise}
\begin{exercise}[$\star\star$]\label{ex:v9-four-dimensional-order-five}
\emph{Fifth-order lattice symmetry in four dimensions.} In $\R^5$, let $W=\{x:\sum_ix_i=0\}$ and $\Lambda=W\cap\Z^5$. (a) Prove that $f_i=e_i-e_5$, $1\le i\le4$, is a lattice basis in $W$. (b) Let $Ce_i=e_{i+1}$, with indices modulo five. Prove that this is an orthogonal transformation of order five on $W$ preserving $\Lambda$, and find its integer matrix and trace in the basis $f_i$. (c) Compute the Gram matrix of this basis and explain why the planar trace restriction does not apply. Identify the two real rotation planes using the cosine and sine vectors of fifth roots of unity.
\end{exercise}

\section{Hints}

\begin{hint}{geometry:ex:chasles}
For intersecting mirrors, put their intersection at zero and the first mirror on the real axis. For parallel mirrors, use $\operatorname{Im}z=0,d$.
\end{hint}
\begin{hint}{geometry:ex:tworotations}
Use $R_{p,\alpha}(z)=e^{i\alpha}z+(1-e^{i\alpha})p$. The composite has linear factor $e^{i(\alpha+\beta)}$; its constant term determines either the fixed point or the translation.
\end{hint}
\begin{hint}{geometry:ex:dihedral}
A symmetry sends the polygon's centre to itself and sends a chosen vertex to one of $n$ vertices. After that, at most two choices remain for its neighbour.
\end{hint}
\begin{hint}{geometry:ex:conjugacy}
For $h(x)=Ax+b$, conjugation moves a rotation centre or mirror axis by $h$ and a translation vector by $A$. An indirect coordinate change reverses rotation angle.
\end{hint}
\begin{hint}{geometry:ex:affineellipse}
A circle chord parallel to $u$ has midpoint on the diameter perpendicular to $u$. Only parallelism and midpoint relations need survive the affine map.
\end{hint}
\begin{hint}{geometry:ex:friezeid}
Use the individual glyph symmetries for the first five rows. In each alternating row, a horizontal reflection exchanges the two glyphs, so combine it with a one-cell shift.
\end{hint}
\begin{hint}{geometry:ex:fivefold}
After replacing the rotation by a suitable power, consider $v+\rho^2v$. Its two summands make angle $144^\circ$.
\end{hint}
\begin{hint}{geometry:ex:hexagonal}
Reduce the real coefficients of a lattice vector in the basis $v,\rho v$ to $[-1/2,1/2]$. The remaining vector has length at most $\sqrt3|v|/2$.
\end{hint}
\begin{hint}{geometry:ex:cube}
A rotation fixing a face and one marked edge fixes two independent vectors from the cube centre. For the axis list, distinguish opposite faces, opposite vertices and opposite edge midpoints.
\end{hint}
\begin{hint}{geometry:ex:billiards}
Fold each real coordinate modulo two into $[0,1]$. A return with unchanged nonzero velocity components requires even reflection parity in those coordinates.
\end{hint}
\begin{hint}{geometry:ex:glidesquare}
Write the glide as $(x,y)\mapsto(x+t,u-y)$ with $t\ne0$. Its square translates by $(2t,0)$.
\end{hint}
\begin{hint}{geometry:ex:finiteSO3}
A pole orbit has $N/n_i$ elements, each fixed by $n_i-1$ nonidentity rotations. Each contribution $1-1/n_i$ lies in $[1/2,1)$.
\end{hint}
\begin{hint}{geometry:ex:ceva}
Use barycentric weights $(\alpha,\beta,\gamma)$. Their sum being nonzero gives a point; sum zero gives a common direction. A finite nonvertex side point has signed ratio different from zero and $-1$.
\end{hint}
\begin{hint}{geometry:ex:cevaparallel}
The vertices must have distinct horizontal coordinates. Read each ratio from the horizontal projections of its side.
\end{hint}
\begin{hint}{geometry:ex:eulerline}
The conjugate of $(b+c)/(b-c)$ is its negative. For coincident centres, a median through the circumcentre is also perpendicular to its opposite chord.
\end{hint}
\begin{hint}{geometry:ex:routh}
Normalize to $(0,0),(1,0),(0,1)$, remembering that its area is $1/2$, not one. Subtract the same intersection point when forming both edge vectors.
\end{hint}
\begin{hint}{ex:v9-finite-affine-group}
Average an orbit for $c$. Right multiplication permutes the summands in $Q$. For $Q=\bigl(\begin{smallmatrix}u&v\\v&w\end{smallmatrix}\bigr)$, use an upper triangular $S$ with entries $\sqrt u,v/\sqrt u,\sqrt{w-v^2/u}$.
\end{hint}
\begin{hint}{ex:v9-glide-wallpaper}
Use $g^2=a$ and $gbg^{-1}=b^{-1}$ to obtain $a^pb^qg^\varepsilon$, with $\varepsilon=0,1$. Solve the affine fixed-point equations for the odd case.
\end{hint}
\begin{hint}{ex:v9-ellipse-parallelogram}
Pull the parallelogram back to the circle. A cyclic parallelogram is a rectangle, and its area is largest when it is a square. Transform every equality case back, scaling areas by $|\det T|$.
\end{hint}
\begin{hint}{ex:v9-four-dimensional-order-five}
The fifth coordinate is minus the sum of the first four. For the cyclic map, $Cf_i=f_{i+1}-f_1$ for $i<4$ and $Cf_4=-f_1$. Split the nonconstant discrete Fourier modes into their real and imaginary parts.
\end{hint}

\section{Solutions}

\begin{solution}{geometry:ex:chasles}
\idea Choose coordinates in which one mirror is simple, then compose the two explicit reflection formulas.
\writeup For intersecting lines put their intersection at zero and $m_1$ on the real axis. Reflection in $m_2$ at directed angle $\varphi$ is $e^{2i\varphi}\bar z$, while the first reflection is $\bar z$. Their composite is $e^{2i\varphi}z$. The angle is understood modulo $2\pi$, since changing a line angle by $\pi$ changes nothing. For parallel mirrors use $\operatorname{Im}z=0$ and $d$. The formulas $s_1(z)=\bar z$, $s_2(z)=\bar z+2id$ give $s_2s_1(z)=z+2id$, twice the signed perpendicular displacement. When $d=0$ it is the identity.
\end{solution}

\begin{solution}{geometry:ex:tworotations}
\idea The coefficient of $z$ records the total angle; the constant term records the remaining displacement.
\writeup Write $R_{p,\alpha}(z)=e^{i\alpha}z+(1-e^{i\alpha})p$. Composing in the stated order gives
\[
R_{q,\beta}R_{p,\alpha}(z)=e^{i(\alpha+\beta)}z+B,
\quad B=e^{i\beta}(1-e^{i\alpha})p+(1-e^{i\beta})q.
\]
If the linear factor is not one, the centre is $B/(1-e^{i(\alpha+\beta)})$ and the angle is $\alpha+\beta$. Otherwise the map is translation by $B$, possibly zero. These formulas include $p=q$: they give rotation about that common point, or the identity for total angle zero. For distinct centres, writing the two rotations as pairs of reflections with a common middle mirror along $pq$ also shows the result geometrically; the other mirrors either intersect at the new centre or are parallel.
\end{solution}

\begin{solution}{geometry:ex:dihedral}
\idea A vertex image and an orientation determine the whole regular polygon.
\writeup Every symmetry fixes the centroid, hence is a rotation or reflection about that centre. A chosen vertex has $n$ possible images, and its adjacent vertex must go to one of the two neighbours of that image, giving at most $2n$ maps. The $n$ rotations by multiples of $2\pi/n$ and their compositions with one polygon mirror realize all $2n$. Orientation gives a homomorphism to $\{1,-1\}$ whose kernel is the rotation subgroup, so that subgroup is normal. With a chosen mirror as the real axis, $s(z)=\bar z$ and $r(z)=e^{i\theta}z$, giving
$srs(z)=\overline{e^{i\theta}\bar z}=e^{-i\theta}z=r^{-1}(z)$.
\end{solution}

\begin{solution}{geometry:ex:conjugacy}
\idea Conjugation moves the geometric data by an isometry, which keeps lengths but may reverse orientation.
\writeup Under $h(x)=Ax+b$, a rotation centre $c$ moves to $h(c)$; its angle stays $\theta$ for direct $h$ and becomes $-\theta$ for indirect $h$. This follows from conjugation by complex rotations/translations and by $z\mapsto\bar z$. Thus equality up to sign is necessary, and translating centres together, with a reflection if needed, proves sufficiency. Angle zero gives the identity independently of its nominal centre.

A translation vector $v$ becomes $Av$, so its length is invariant. Any equal-length pair, including zero, is related by an orthogonal matrix, proving sufficiency. A mirror line becomes $h(\ell)$, and any line can be moved to any other, so all reflections are conjugate. A glide's axis and displacement $v$ become $h(\ell),Av$. The length is necessary, also detectable from its square's translation length $2|v|$. An isometry mapping an oriented first axis to the second with the prescribed equal-length displacement conjugates the glides, proving sufficiency.
\end{solution}

\begin{solution}{geometry:ex:affineellipse}
\idea In the circle, equal endpoint radii put a chord's midpoint on a perpendicular diameter. The affine map preserves the midpoint and line, even though it need not preserve their right angle.
\writeup For endpoints $x,y$ on a centred circle, their midpoint $m$ satisfies
$\langle m,y-x\rangle=(|y|^2-|x|^2)/2=0$. Thus chords parallel to $u$ have midpoints on the diameter perpendicular to $u$. Equivalently chords perpendicular to $u$ have midpoints on the diameter parallel to $u$. If $f$ sends the circle to the ellipse, inverse images of a parallel family remain parallel, $f(m)=(f(x)+f(y))/2$, and the midpoint diameter maps to a line through the image centre. This proves the claim; tangent limiting chords obey it by continuity as well.
\end{solution}

\begin{solution}{geometry:ex:friezeid}
\idea List the axis and side signs that carry one ideal cell to a suitable cell. Alternation may require a shift before the reflected glyph matches.
\writeup Let the cell spacing be $d$. For the first five rows the least translation is $d$. The ideal $L$ has no additional symmetry, giving a hop. The ideal $N$ has half-turns about cell centres and no mirrors, giving a spinning hop. Vertical mirrors through and between the $V$ cells give a sidle; horizontal reflection would produce absent $\Lambda$ cells. The $D$ row has the horizontal mirror but no vertical mirror or half-turn, giving a jump. The $H$ row has both mirrors and their half-turn products, giving a spinning jump.

For the $pb$ row the least translation is $2d$. Horizontal reflection exchanges the two asymmetric glyphs, and a shift by $d$ places each over the other, giving the half-period glide. No vertical mirror or half-turn can match their asymmetric stroke arrangements. This is a step. In the alternating $V\Lambda$ row the least translation is also $2d$, with the same half-period horizontal glide. Vertical mirrors through glyph centres also preserve the row; their products with glides give half-turns between adjacent cells. There is no horizontal mirror alone. This is a spinning sidle. These statements use the ideal shapes and separations stipulated in the exercise; a font with unequal bowls or stems need not have those exact symmetries.
\end{solution}

\begin{solution}{geometry:ex:fivefold}
\idea Two rotated shortest vectors can add to a vector shorter than either. A lattice cannot contain that smaller nonzero sum.
\writeup The linear part of any lattice-preserving affine rotation also preserves the lattice, as proved in the restriction theorem. Replace a fifth-order linear rotation by an appropriate power so its angle is $72^\circ$. For shortest nonzero $v$, $\rho^2v$ has the same length $\delta$ and makes angle $144^\circ$ with it. Their sum is a lattice vector of length $2\delta\cos72^\circ$. This is positive because $72^\circ<90^\circ$, and less than $\delta$ because $72^\circ>60^\circ$. The contradiction requires no trace or decimal approximation.
\end{solution}

\begin{solution}{geometry:ex:hexagonal}
\idea A missing lattice vector would leave a short nonzero remainder after subtracting the obvious two generators.
\writeup The real vectors $v,\rho v$ form a basis with angle $60^\circ$. For $x\in\Lambda$, subtract integer multiples to leave $av+b\rho v$ with $|a|,|b|\le1/2$. Its squared norm is
$|v|^2(a^2+b^2+ab)\le3|v|^2/4$. A nonzero such remainder contradicts minimality, so it is zero and the two vectors generate $\Lambda$. In complex coordinates $\rho v=e^{i\pi/3}v=(1+\omega)v$. Since $\omega^2+\omega+1=0$, the ring $\Z[\omega]$ is $\Z+\Z\omega=\Z+\Z e^{i\pi/3}$. Multiplication by the nonzero complex number $v$ is the desired similarity.
\end{solution}

\begin{solution}{geometry:ex:cube}
\idea A marked face fixes an oriented frame. The possible frames count the group; an axis inventory provides a second check.
\writeup Place the cube centrally with its faces perpendicular to the coordinate axes. A face with a marked edge determines its outward unit normal and the unit vector from face centre to edge midpoint. These are perpendicular, and their cross product completes an oriented orthonormal frame. There are $6\cdot4=24$ possible target flags. Mapping one frame to another gives a unique signed coordinate permutation of determinant one, which preserves the cube. Thus all flags occur and the stabilizer is trivial, giving $24$ rotations.

For the axis count: there is the identity; $3\cdot3=9$ nonidentity rotations by $90^\circ,180^\circ,270^\circ$ about opposite-face axes; $4\cdot2=8$ rotations by $120^\circ,240^\circ$ about opposite-vertex axes; and six half-turns about opposite-edge-midpoint axes. These maps are distinct by their axes and angles and total $24$, so the flag bound shows that none are missing.

Their action on four body-diagonal lines identifies the group with $S_4$. To check faithfulness, represent those lines by $(1,1,1),(1,-1,-1),(-1,1,-1),(-1,-1,1)$. Their mutual inner products are $-1$. A rotation fixing all four lines sends each representative to itself times a sign; preserving their inner products forces all signs equal. They span space, so the map is $I$ or $-I$, and the latter has determinant $-1$. The kernel is trivial; both groups have order $24$, giving the identification.
\end{solution}

\begin{solution}{geometry:ex:billiards}
\idea Reflect the square at a collision and continue straight. Folding a straight line back retains the position and reflection parity.
\writeup Begin at an interior point of the nonsingular trajectory. Unfolded squares tile the plane; folding each coordinate modulo two, with $s$ in $[1,2]$ sent to $2-s$, maps the straight path $x_0+t(q,p)$ back to the table. After $t=2$ its displacement is $(2q,2p)$, so both position and velocity return. The travelled length is $2\sqrt{q^2+p^2}$.

For any earlier positive return with unchanged direction, a coordinate with nonzero velocity must have the same folding parity as initially. Since the initial point is interior, that requires its displacement to lie in $2\Z$. A zero velocity coordinate has zero displacement anyway. Thus $tq,tp\in2\Z$. Bezout's identity for the coprime pair implies $t/2\in\Z$, so the least positive return has $t=2$. For $p=0,q=1$ the length is two. A vertical trajectory similarly has length two, using direction $(0,1)$. Corner hits are excluded because the two side reflections there do not specify a single ordinary billiard collision; the unfolded formula alone does not supply a physical corner rule.
\end{solution}

\begin{solution}{geometry:ex:glidesquare}
\idea Squaring removes the reflection and doubles the along-axis displacement. It cannot double a nonzero displacement to zero.
\writeup Choose coordinates parallel to the axis. A glide is $g(z)=\bar z+t+iu$ with real $t\ne0$, whose square is $z\mapsto z+2t$. The axis is $\operatorname{Im}z=u/2$. Conversely, for translation by a fixed nonzero vector $v$, a glide with that square must have displacement $v/2$ and an axis parallel to $v$. Reflection in any such axis followed by $v/2$ has the required square, and the chosen axis determines it uniquely. The identity translation is not the square of a genuine glide. It is the square of any reflection, which has zero along-axis displacement and belongs to a different class.
\end{solution}

\begin{solution}{geometry:ex:finiteSO3}
\idea Count the same finite set once by rotations and once by pole orbits. The resulting reciprocal sum has very few possibilities.
\writeup Each of the $N-1$ nonidentity rotations fixes exactly two poles, giving $2(N-1)$ pairs. A pole with stabilizer order $n_i\ge2$ belongs to an orbit of size $N/n_i$, and each point of that orbit contributes $n_i-1$ pairs. Equating counts and dividing by $N$ gives the stated formula.

Its left side lies in $[1,2)$ and each summand in $[1/2,1)$, so there must be two or three orbits. For two, $1/n_1+1/n_2=2/N$. Since each $n_i\le N$, both must equal $N$, giving $(n,n)$ with $N=n\ge2$. For three, sort $2\le n_1\le n_2\le n_3$. Their reciprocal sum is $1+2/N>1$, forcing $n_1=2$. Then $1/n_2+1/n_3>1/2$ forces $n_2=2$, or $n_2=3$ and $n_3=3,4,5$. Substitution yields $N=2n$ for $(2,2,n)$, and $12,24,60$ for the last three lists. The counting gives the required stabilizer lists and orders; the familiar dihedral and Platonic rotation examples realize them.
\end{solution}

\begin{solution}{geometry:ex:ceva}
\idea Barycentric weights encode the three line conditions. A nonzero sum of weights gives their meeting point; a zero sum gives their common direction.
\writeup Every finite point has weights $(\alpha,\beta,\gamma)$ with sum nonzero, unique up to common scaling, through
$P=(\alpha A+\beta B+\gamma C)/(\alpha+\beta+\gamma)$. This follows by expressing $P-A$ in the basis $B-A,C-A$. If the three cevians meet at $P$, no weight is zero: for example $\alpha=0$ would put $P$ on $BC$, making $AD$ meet $BC$ at $D=P$ and then forcing $BE=BC$, contrary to the exclusion of side vertices. None of the pair sums is zero either. For instance $\beta+\gamma=0$ would make $P-A$ parallel to $B-C$, so the line $AP$ could not meet $BC$ at a finite $D$. Therefore
\[
D=\frac{\beta B+\gamma C}{\beta+\gamma},\quad
E=\frac{\gamma C+\alpha A}{\gamma+\alpha},\quad
F=\frac{\alpha A+\beta B}{\alpha+\beta}.
\]
Their signed ratios are $\gamma/\beta,\alpha/\gamma,\beta/\alpha$, whose product is one.

For the parallel case choose their common direction as vertical. The vertex horizontal coordinates are distinct: if two agreed, their side would be vertical and the third vertex's parallel cevian could not meet that side finitely. Since $D,E,F$ have horizontal coordinates $x_A,x_B,x_C$, respectively, the product is
\[
\frac{x_A-x_B}{x_C-x_A}\,
\frac{x_B-x_C}{x_A-x_B}\,
\frac{x_C-x_A}{x_B-x_C}=1.
\]
Thus both possible configurations imply the condition.

Conversely suppose the product is one. Put $\alpha=1$, $\beta=AF/FB$ and $\gamma=EA/CE$. All are nonzero and the product condition gives $\gamma/\beta=BD/DC$. Each pair sum is nonzero: a finite side point never has signed ratio $-1$, because $D-B=-(C-D)$ would force $B=C$, and the other sides are the same. If $S=\alpha+\beta+\gamma\ne0$, the finite point with these weights lies on all three cevians by the displayed formulas. If $S=0$, set $V=\alpha A+\beta B+\gamma C$. It is nonzero by affine independence. For example
$D-A=V/(\beta+\gamma)$, and the cyclic expressions give the same common direction for $E-B,F-C$. Hence the cevians are parallel. This proves the converse in full without treating parallelism as a finite point.
\end{solution}

\begin{solution}{geometry:ex:cevaparallel}
\idea Project along the common direction onto a transverse coordinate axis. Signed ratios on each side are preserved by this one-dimensional affine projection.
\writeup Choose the cevians vertical. No side is vertical, since a parallel line through its opposite vertex would then miss it. Thus $x_A,x_B,x_C$ are pairwise distinct. Reading the three side ratios at the horizontal coordinates of $D,E,F$ gives
\[
\frac{BD}{DC}=\frac{x_A-x_B}{x_C-x_A},\quad
\frac{CE}{EA}=\frac{x_B-x_C}{x_A-x_B},\quad
\frac{AF}{FB}=\frac{x_C-x_A}{x_B-x_C}.
\]
Every denominator is nonzero and the product cancels to one. The cancellation holds for signed ratios, so points on the extensions of the sides require no separate case.
\end{solution}

\begin{solution}{geometry:ex:eulerline}
\idea The unit-circle condition makes conjugation into reciprocation. That turns the altitude equations into a simple sum of the vertex vectors.
\writeup For $h=a+b+c$, the altitude direction from $a$ is $h-a=b+c$. Since $b\ne c$ and $\bar b=1/b$, $\bar c=1/c$,
\[
\overline{\frac{b+c}{b-c}}=
\frac{1/b+1/c}{1/b-1/c}=-\frac{b+c}{b-c}.
\]
The quotient is purely imaginary, so $h-a$ is perpendicular to side $bc$; a zero vector simply means $h=a$, as can happen in a right triangle. The other altitude conditions follow cyclically, proving that $h$ is the orthocentre. The circumcentre is zero and $h=3g$, so $|h-g|=2|g|$, and for $g\ne0$ the centroid is strictly between the other two centres.

If $g=0$, the median from each vertex passes through the circumcentre and is perpendicular to its opposite chord. It is therefore also that side's perpendicular bisector, forcing the two adjacent side lengths to be equal. Applying this at the vertices makes the triangle equilateral. Conversely an equilateral triangle has centroid, circumcentre and orthocentre equal. This includes the coincident-centre case without requiring a unique line through three equal points.
\end{solution}

\begin{solution}{geometry:ex:routh}
\idea Normalize the triangle affinely, compute one small determinant, then take an area ratio so that the normalization disappears.
\writeup Use $A=(0,0)$, $B=(1,0)$, $C=(0,1)$, whose area is $1/2$. The consistent division points are $D=(2/3,1/3)$, $E=(0,2/3)$ and $F=(1/3,0)$. Solving the pairs of line equations gives
\[
P=AD\cap BE=(4/7,2/7),\quad
Q=BE\cap CF=(1/7,4/7),\quad
R=CF\cap AD=(2/7,1/7).
\]
Both edge vectors must start at $P$:
\[
Q-P=(-3/7,2/7),\qquad R-P=(-2/7,-1/7).
\]
Their determinant is $3/49+4/49=1/7$, so the small area is $1/14$, exactly one seventh of $1/2$. An invertible affine map multiplies both areas by the same positive determinant magnitude; in a triangle of area one the enclosed area is therefore $1/7$.
\end{solution}

\begin{solution}{ex:v9-finite-affine-group}
\idea Average both a point and a squared-length form over the finite group. The point average fixes an origin; the form average fixes a ruler.
\writeup (a) The orbit average $c=|G|^{-1}\sum_g g(p)$ is fixed because affine maps preserve averages and left multiplication permutes the group. After moving $c$ to zero every map is linear. For nonzero $v$,
$v^\top Qv=\sum_g|A_gv|^2>0$, since the identity term is positive. Moreover
\[
A_h^\top QA_h=\sum_g(A_gA_h)^\top(A_gA_h)=Q,
\]
by the permutation $g\mapsto gh$.

(b) Write $Q=\bigl(\begin{smallmatrix}u&v\\v&w\end{smallmatrix}\bigr)$. Positivity gives $u>0$ and $w-v^2/u>0$, the latter by testing $(-v/u,1)$. Thus
\[
S=\begin{pmatrix}\sqrt u&v/\sqrt u\\0&\sqrt{w-v^2/u}\end{pmatrix}
\]
is invertible and direct multiplication gives $S^\top S=Q$. The conjugated matrices $B_h=SA_hS^{-1}$ obey
$B_h^\top B_h=S^{-\top}A_h^\top QA_hS^{-1}=I$. Leonardo's theorem classifies the resulting finite orthogonal group as cyclic or dihedral, including the trivial group. This is an affine coordinate change, which need not preserve the original Euclidean lengths.

(c) The matrix $A=\bigl(\begin{smallmatrix}0&2\\1/2&0\end{smallmatrix}\bigr)$ has $A^2=I$, so $\{I,A\}$ is a finite affine group. But $|Ae_1|=1/2\ne1$, so it is not a Euclidean isometry group in these coordinates. Its invariant form from (a) is $Q=\operatorname{diag}(5/4,5)$; rescaling by the corresponding $S$ makes $A$ the ordinary coordinate-swap reflection.
\end{solution}

\begin{solution}{ex:v9-glide-wallpaper}
\idea Separate integer translations from the one remaining reflection sign. That sign forces a half-integer horizontal displacement, which prevents any mirror from fixing points.
\writeup (a) We have $a(x,y)=(x+1,y)$, $gbg^{-1}=b^{-1}$, and $a$ commutes with both generators. Therefore every word reduces to
\[
a^pb^qg^\varepsilon(x,y)=
(x+p+\varepsilon/2,(-1)^\varepsilon y+q),
\quad p,q\in\Z,\quad\varepsilon\in\{0,1\}.
\]
Conversely all these maps are group elements. Equality of two formulas first fixes the sign $\varepsilon$, then fixes $p,q$, proving uniqueness. The translations are precisely the even-sign maps, giving lattice $\Z^2$ and quotient $\Z/2$.

(b) A nonidentity even-sign map is a nonzero translation and fixes nothing. An odd-sign map has horizontal displacement $p+1/2\ne0$, so also fixes nothing. It reflects in $y=q/2$ and translates along that axis by $p+1/2$. Thus these and only these are its glide axes and displacements; in particular it contains no reflection.

(c) A power of $g$ reduces the horizontal coordinate uniquely into $[0,1/2)$; its possible sign flip does not obstruct then reducing the vertical coordinate modulo one using $b$. This gives an orbit representative in the stated half-open rectangle. If two representatives there were related, the horizontal displacement, an integer multiple of $1/2$, would have to vanish. Uniqueness of the normal form then forces $\varepsilon=p=0$, and the vertical interval forces $q=0$. The representative is unique, also on the included lower/left boundaries. The comparison group contains a genuine mirror fixing a line. Conjugation preserves the existence of nonidentity elements with fixed points, so it cannot be conjugate to $G$, even by a general homeomorphism.
\end{solution}

\begin{solution}{ex:v9-ellipse-parallelogram}
\idea Circle geometry gives both the maximum and every equality case. An affine map changes area by one factor, so none of those maximizers is lost.
\writeup (a) Apply $T^{-1}$ after subtracting $c$. An inscribed parallelogram becomes a cyclic parallelogram in $S^1$. Opposite angles of a parallelogram are equal, while opposite angles of a cyclic quadrilateral are supplementary, so all are right angles. Its centre is the circle centre. Write its side half-vectors as $a,b$, with $a\perp b$ and $|a|^2+|b|^2=1$. Its area is
$4|a||b|\le2(|a|^2+|b|^2)=2$, with equality exactly when $|a|=|b|=1/\sqrt2$, a square. Degenerate parallelograms have area zero and cannot maximize.

Every affine image has area multiplied by $|\det T|$, so the maximum is $2|\det T|$. All equality cases are images of inscribed squares, equivalently parallelograms with vertices
\[
c+Tu,\quad c+Tv,\quad c-Tu,\quad c-Tv,
\qquad |u|=|v|=1,\quad u\perp v.
\]
Here $u,v$ are the square's unit half-diagonals, giving a family of maximizers through all orientations.

(b) For $T=\operatorname{diag}(2,1)$, $|\det T|=2$ and the maximum is four. The diamond has perpendicular diagonal lengths four and two, hence area four. The rectangle has side lengths $2\sqrt2$ and $\sqrt2$, again area four. The diamond's adjacent side vectors, for example $(-2,1)$ and $(-2,-1)$ at $(2,0)$, have dot product three, so it is not a Euclidean rectangle. Affine maps preserve the area ratio and parallelogram relations, while they may change right angles.
\end{solution}

\begin{solution}{ex:v9-four-dimensional-order-five}
\idea A fifth-order orthogonal symmetry is possible when two rotation planes share the trace. The lattice basis need not be orthonormal.
\writeup (a) Every $x\in W$ is $\sum_{i=1}^4x_i(e_i-e_5)$, with unique first four coordinates. Integer $x$ therefore correspond exactly to integer coefficients in the basis $f_i$. Thus $\Lambda$ is a full-rank lattice in the four-dimensional Euclidean space $W$.

(b) Take $Ce_i=e_{i+1}$, with indices modulo five. Coordinate permutation is orthogonal, preserves integer coordinates and their sum, and has $C^5=I$. On $W$ it is not the identity, so its order, a nontrivial divisor of five, is five. From $Cf_i=f_{i+1}-f_1$ for $i<4$ and $Cf_4=-f_1$, its matrix is
\[
A=\begin{pmatrix}-1&-1&-1&-1\\1&0&0&0\\0&1&0&0\\0&0&1&0\end{pmatrix},
\qquad\operatorname{tr}A=-1.
\]

(c) The basis inner products are $\langle f_i,f_j\rangle=\delta_{ij}+1$, so its Gram matrix is $Q=I_4+J_4$, where $J_4$ has every entry one. Orthogonality is $A^\top QA=Q$ in this basis, which also follows directly by multiplying. It does not require $A^\top A=I$ for this nonorthonormal basis.

For $j=1,2$, form vectors $c_j=(\cos(2\pi jk/5))_{k=0}^4$ and $s_j=(\sin(2\pi jk/5))_{k=0}^4$. Finite geometric sums of fifth roots give coordinate sums zero, squared lengths $5/2$, and mutual orthogonality of these four vectors: apply the sum to frequencies $j$, $2j$, and $j\pm l$. They form an orthogonal basis of $W$. The cyclic shift maps $c_j$ to $\cos\theta_j c_j+\sin\theta_j s_j$ and $s_j$ to $-\sin\theta_j c_j+\cos\theta_j s_j$, where $\theta_j=2\pi j/5$. Thus it rotates two perpendicular planes through $72^\circ$ and $144^\circ$. Their trace contributions add to
$2\cos72^\circ+2\cos144^\circ=-1$. The planar argument would require a single $2\cos\theta$ to be integral, and the three-dimensional argument only adds a fixed-axis contribution one. Neither describes this four-dimensional example.
\end{solution}

\section*{Sources and further reading}

\begin{itemize}
\item F.~Klein, \emph{Vergleichende Betrachtungen \"uber neuere geometrische Forschungen}, Erlangen, 1872; English translation \emph{A
comparative review of recent researches in geometry}, Bulletin of the New York Mathematical Society 2 (1893), 215--249.
\item M.~A.~Armstrong, \emph{Groups and Symmetry}, Springer, 1988: a route from plane isometries to frieze and wallpaper classifications, followed by finite rotation groups.
\item J.~H.~Conway, H.~Burgiel and C.~Goodman-Strauss, \emph{The Symmetries of Things}, A~K~Peters, 2008: for the frieze names and the orbifold description of pattern groups.
\item H.~Weyl, \emph{Symmetry}, Princeton University Press, 1952: read its treatment of ornaments and crystals alongside the finite-group arguments here.
\item G.~P\'olya, \emph{\"Uber die Analogie der Kristallsymmetrie in der Ebene}, Zeitschrift f\"ur Kristallographie 60 (1924), 278--282;
D.~Schattschneider, \emph{The plane symmetry groups: their recognition and notation}, American Mathematical Monthly 85 (1978), 439--450.
\item D.~Shechtman, I.~Blech, D.~Gratias and J.~W.~Cahn, \emph{Metallic phase with long-range orientational order and no translational
symmetry}, Physical Review Letters 53 (1984), 1951--1953.
\item T.~Hawkins, \emph{The Erlanger Programm of Felix Klein: reflections on its place in the history of mathematics}, Historia
Mathematica 11 (1984), 442--470 (Klein's appointment and the reception of the programme).
\item D.~Schattschneider, \emph{Visions of Symmetry: Notebooks, Periodic Drawings, and Related Work of M.~C.~Escher}, W.~H.~Freeman, 1990
(Escher's visits to the Alhambra in 1922 and 1936 and his reading of P\'olya).
\item D.~K.~Washburn and D.~W.~Crowe, \emph{Symmetries of Culture: Theory and Practice of Plane Pattern Analysis}, University of Washington
Press, 1988 (frieze types in anthropology).
\item T.~L.~Heath, \emph{A History of Greek Mathematics}, vol.~2, Oxford, 1921 (Heron's proof of the law of reflection in his
\emph{Catoptrica}).
\item H.~S.~M.~Coxeter, \emph{Introduction to Geometry}, 2nd ed., Wiley, 1969, Chapters~3 and~4 (isometries, Fagnano's problem).
\end{itemize}
